\documentclass{article}
\usepackage[T1]{fontenc}
\usepackage{lmodern}
\usepackage{graphicx}
\usepackage{amsmath,amssymb}
\usepackage[a4paper,margin=2cm]{geometry}
\usepackage[absolute,overlay]{textpos}
\setlength{\TPHorizModule}{1mm}
\setlength{\TPVertModule}{1mm}
\usepackage[indonesian]{babel}
\usepackage{hyperref}
\setlength{\emergencystretch}{2em}
% unit: Halaman 70

\begin{document}

% === Halaman 70 (llm) ===

\section*{Pembangkitan kunci putaran}

\begin{itemize}
    \item Setiap putaran di dalam \textit{iterated cipher} menggunakan kunci internal (\textit{subkey}) yang dinamakan kunci putaran (\textit{round key}).
    \item Kunci putaran dibangkitkan dari kunci eksternal yang diberikan oleh \textit{user} melalui proses yang dinamakan \textit{key expansion} atau \textit{key scheduling}.
    \item Di dalam \textit{key expansion} dilakukan komputasi yang kompleks untuk menghasilkan sejumlah kunci putaran yang berbeda-beda.
    \item Panjang kunci putaran tidak selalu sama dengan panjang kunci eksternal.
\end{itemize}

\end{document}
